\subsection{Создание, изменение и удаление томов}

Том GlusterFS собирается из бриков~--- каталогов на серверных узлах. На \code{node1}, \code{node2} и \code{node3} созданы каталоги бриков (рисунки~\ref{fig:dir1}--\ref{fig:dir3}):
\begin{Verbatim}
sudo mkdir -p /var/brick1 /var/brick2 /var/brick3 /var/brick4 /var/brick5
\end{Verbatim}
Каталоги в \code{/var} (на корневом разделе) используются в учебных целях, поэтому при создании томов применяется ключ \code{force}. В реальных системах брики размещают на отдельных дисках.

\fig{4.1. Каталоги бриков(node1).png}{Каталоги бриков на node1}{fig:dir1}
\fig{4.1. Каталоги бриков(node2).png}{Каталоги бриков на node2}{fig:dir2}
\fig{4.1. Каталоги бриков(node3).png}{Каталоги бриков на node3}{fig:dir3}

\subsubsection{Тома replica, replica с арбитром и distributed}

На узле \code{node1} созданы три тома (рисунок~\ref{fig:create}):
\begin{itemize}
  \item \code{volume1}~--- том типа replica с двумя копиями (\code{node1}, \code{node2}). Каждый файл хранится на обоих бриках. При создании GlusterFS предупреждает, что replica~2 подвержена split-brain (расхождению копий при потере связи между узлами), и просит подтверждения;
  \item \code{vol_arb}~--- том replica~3 с арбитром (\code{replica 3 arbiter 1}). Третий брик (\code{node3}) хранит только метаданные файлов и не требует места под данные, но позволяет принимать решение о кворуме и защищает от split-brain;
  \item \code{volume2}~--- том типа distributed, в котором файлы распределяются между бриками целиком (каждый файл лежит на одном из бриков).
\end{itemize}

\begin{Verbatim}
sudo gluster volume create volume1 replica 2 node1:/var/brick1 node2:/var/brick1 force
sudo gluster volume start volume1

sudo gluster volume create vol_arb replica 3 arbiter 1 \
  node1:/var/brick3 node2:/var/brick3 node3:/var/brick3 force
sudo gluster volume start vol_arb

sudo gluster volume create volume2 node1:/var/brick2 node2:/var/brick2 force
sudo gluster volume start volume2
\end{Verbatim}

\fig{4.2. Создание томов (всё на node1).png}{Создание и запуск томов volume1, vol\_arb и volume2}{fig:create}

Тома смонтированы на клиенте с помощью файловой системы \code{glusterfs} (рисунок~\ref{fig:mount}). В тома записаны тестовые файлы, затем проверено их размещение на бриках:
\begin{Verbatim}
sudo mount -t glusterfs node1:/volume1 /mnt/volume1
sudo mount -t glusterfs node1:/volume2 /mnt/volume2
sudo mount -t glusterfs node1:/vol_arb /mnt/vol_arb
\end{Verbatim}

\fig{4.3. Монтирование и проверка на client.png}{Монтирование томов и запись тестовых файлов на клиенте}{fig:mount}

\textbf{Replica.} Файл \code{test0.txt}, записанный в \code{/mnt/volume1}, присутствует на бриках обоих узлов (рисунки~\ref{fig:b1n1}, \ref{fig:b1n2}).

\fig{4.3. ls -l -var-brick1(node1).png}{Содержимое /var/brick1 на node1}{fig:b1n1}
\fig{4.3. ls -l -var-brick1(node2).png}{Содержимое /var/brick1 на node2}{fig:b1n2}

\textbf{Replica с арбитром.} Файл, записанный в \code{/mnt/vol_arb}, присутствует на всех трёх узлах, однако на узле-арбитре \code{node3} его размер равен нулю: арбитр хранит лишь метаданные (рисунки~\ref{fig:b3n1}--\ref{fig:b3n3}).

\fig{4.3. ls -l -var-brick3(node1).png}{Содержимое /var/brick3 на node1}{fig:b3n1}
\fig{4.3. ls -l -var-brick3(node2).png}{Содержимое /var/brick3 на node2}{fig:b3n2}
\fig{4.3. ls -l -var-brick3(node3).png}{Содержимое /var/brick3 на node3 (арбитр)}{fig:b3n3}

\textbf{Distributed.} В \code{/mnt/volume2} создано 10 файлов. Они поделены между бриками \code{node1} и \code{node2}: каждый файл целиком хранится на одном из них (рисунки~\ref{fig:b2n1}, \ref{fig:b2n2}).

\fig{4.3. ls -l -var-brick2(node1).png}{Содержимое /var/brick2 на node1}{fig:b2n1}
\fig{4.3. ls -l -var-brick2(node2).png}{Содержимое /var/brick2 на node2}{fig:b2n2}

\subsubsection{Том stripe}

Тип тома \code{stripe}, разбивавший файлы на фрагменты (полосы) по брикам, был объявлен устаревшим и удалён из GlusterFS начиная с версии~6. В использованной версии~10.3 создать такой том невозможно: попытка создания завершается ошибкой (рисунок~\ref{fig:stripe}).
\begin{Verbatim}
sudo gluster volume create vol_stripe stripe 2 \
  node1:/var/brick4 node2:/var/brick4 force
\end{Verbatim}

\fig{4.4. Stripe(node1).png}{Попытка создания тома типа stripe}{fig:stripe}

Функциональность разбиения файла на части по узлам в современных версиях реализуют два механизма, оба из которых проверены в работе.

\textbf{Аналог 1: disperse (erasure coding).} Файл делится на фрагменты, к которым добавляется избыточность (код стирания). При конфигурации из трёх бриков с \code{redundancy 1} данные хранятся в двух фрагментах, третий хранит чётность, том переживает потерю любого одного узла.
\begin{Verbatim}
sudo gluster volume create vol_disp disperse 3 redundancy 1 \
  node1:/var/brick4 node2:/var/brick4 node3:/var/brick4 force
sudo gluster volume start vol_disp
\end{Verbatim}

\fig{4.4. Аналог 1: disperse (erasure coding) (node1).png}{Создание тома vol\_disp на node1 и размещение файла на брике}{fig:disp1}
\fig{4.4. Аналог 1: disperse (erasure coding)(client).png}{Запись файла в том vol\_disp на клиенте}{fig:disp2}

\textbf{Аналог 2: шардирование (sharding).} Включается параметром тома \code{features.shard}; файлы режутся на блоки заданного размера (4~МБ), которые распределяются между бриками. Параметр включён на пустом томе \code{vol_shard}.
\begin{Verbatim}
sudo gluster volume create vol_shard node1:/var/brick5 node2:/var/brick5 force
sudo gluster volume set vol_shard features.shard on
sudo gluster volume set vol_shard features.shard-block-size 4MB
sudo gluster volume start vol_shard
\end{Verbatim}

Блоки хранятся в служебном каталоге \code{.shard} каждого брика (рисунки~\ref{fig:shardc}--\ref{fig:shard2}).

\fig{4.4. Аналог 2: шардирование (client).png}{Запись файла в том vol\_shard на клиенте}{fig:shardc}
\fig{4.4. Аналог 2: шардирование (node1).png}{Создание тома vol\_shard на node1}{fig:shard0}
\fig{4.4. sudo ls -lh -var-brick5-(node1).shard.png}{Блоки файла в каталоге .shard на node1}{fig:shard1}
\fig{4.4. sudo ls -lh -var-brick5-(node2).shard.png}{Блоки файла в каталоге .shard на node2}{fig:shard2}

\subsubsection{Изменение томов}

Поведение тома изменяется параметрами (опциями) командой \code{gluster volume set}. Для тома \code{volume1} изменён размер кэша и ограничен доступ к тому адресом клиента, затем параметры возвращены к значениям по умолчанию командой \code{volume reset} (рисунок~\ref{fig:set}).
\begin{Verbatim}
sudo gluster volume set volume1 performance.cache-size 256MB
sudo gluster volume set volume1 auth.allow 10.0.2.10
sudo gluster volume info volume1
sudo gluster volume reset volume1 performance.cache-size
sudo gluster volume reset volume1 auth.allow
\end{Verbatim}

\fig{4.5. Изменение томов.png}{Изменение и сброс параметров тома volume1}{fig:set}

Изменение состава бриков (добавление, удаление, смена числа реплик) рассмотрено в разделе~\ref{sec:bricks}.

\subsubsection{Удаление томов}

Временные тома \code{vol_disp} и \code{vol_shard} отмонтированы на клиенте, остановлены и удалены (рисунок~\ref{fig:del}). Остальные тома сохранены для дальнейшей работы. Данные в каталогах бриков при удалении тома остаются, поэтому перед повторным использованием каталога его необходимо очистить.
\begin{Verbatim}
sudo gluster volume stop vol_disp
sudo gluster volume delete vol_disp
sudo gluster volume stop vol_shard
sudo gluster volume delete vol_shard
sudo gluster volume list
\end{Verbatim}

\fig{4.6. Удаление томов.png}{Остановка и удаление томов vol\_disp и vol\_shard}{fig:del}
